\begin{frame}{Final Pipeline}
  \centering
  \begin{tikzpicture}[
    node distance=0.75cm,
    every node/.style={font=\small},
    box/.style={
      rectangle,
      rounded corners,
      draw=accent,
      fill=soft,
      align=center,
      minimum width=0.78\textwidth,
      minimum height=0.72cm
    },
    arr/.style={-{Latex[length=2.5mm]}, very thick, draw=accent}
  ]
    \node[box] (data) {Select sparse posed RGB-D views and scene-level train / validation / test splits};
    \node[box, below=of data] (subset) {Run MV-DUSt3R+ to produce candidate point cloud and confidence};
    \node[box, below=of subset] (base) {Use source depth maps with poses/intrinsics to anchor source rays};
    \node[box, below=of base] (heads) {Apply Run 30 selective source-depth correction};
    \node[box, below=of heads] (eval) {Evaluate overall, occlusion, and repeated/wrong-depth gates};

    \draw[arr] (data) -- (subset);
    \draw[arr] (subset) -- (base);
    \draw[arr] (base) -- (heads);
    \draw[arr] (heads) -- (eval);
  \end{tikzpicture}
\end{frame}

\begin{frame}{Experiment Variables}
  \begin{columns}[T]
    \column{0.48\textwidth}
    \textbf{Controlled factors}
    \begin{itemize}
      \item Number of input views
      \item RGB-D source depth availability
      \item Scene-level train / validation / test split
      \item Fixed Run 30 policy on final test
    \end{itemize}

    \column{0.48\textwidth}
    \textbf{Observed outcomes}
    \begin{itemize}
      \item Reconstruction accuracy
      \item Completeness and consistency
      \item Runtime per scene
      \item Occlusion and repeated/wrong-depth failures
    \end{itemize}
  \end{columns}
\end{frame}

\begin{frame}{RGB-only Diagnostic Path}
  \footnotesize
  \begin{enumerate}
    \item First build a supervised label cache with per-view visibility,
    occlusion, floating/wrong-depth, and match labels.
    \item Mine occlusion-heavy, low-overlap, and hard-negative subsets before
    training the learned heads.
    \item Freeze MV-DUSt3R+ and train small OARH / RSDH heads from generated
    labels, excluding direct target-label leakage features.
    \item Runs 22--26 move the heads into reconstruction-level validation gates
    with fixed-confidence, top-k, and all-candidate baselines.
    \item Run 27 is self-contained and learns a bounded confidence residual
    from actual candidates, self-geometry support, and raw image-patch cues.
    \item Completed Run 27 still selects all candidates: learned filtering beats
    fixed confidence but loses completeness against full retention.
    \item These runs motivate switching the final inference contract to RGB-D.
  \end{enumerate}
\end{frame}

\begin{frame}{RGB-D Source-Depth Correction}
  \footnotesize
  \begin{enumerate}
    \item Run 28 therefore preserves candidate count and learns source-ray
    3D corrections from training depth and poses, using no depth at inference.
    \item Completed Run 28 still fails the gate, but its source-depth oracle
    shows large occlusion and ambiguity headroom.
    \item Run 29 approximates that oracle with RGB-only monocular depth aligned
    through the input poses and intrinsics.
    \item Completed Run 29 regresses both occlusion and ambiguity, so RGB-only
    monodepth is not enough.
    \item Run 30 makes source depth an explicit RGB-D input and gates
    full/selective source-ray correction policies.
    \item Selected policy: \texttt{rgbd\_residual\_ge\_0.30}, residual mode,
    \(\alpha = 1.0\), threshold \(0.30m\).
    \item Final selected method: \texttt{rgbd\_source\_depth\_selected}.
    \item Keep final test clean: no threshold tuning, no checkpoint selection
    on test scenes.
  \end{enumerate}
\end{frame}
